\chapter*{Bibliografía}

\renewcommand{\refname}{Bibliografía}
\begin{thebibliography}{99}

    \bibitem{barraza2021} Barraza, O., Estrada, M. (2021). Structural Analysis in Transit System Using Network Theory: Case of Guadalajara, Mexico. \textit{Urban Sci}, 5, 87. https://doi.org/10.3390/urbansci5040087
    \bibitem{zhao2024} Zhao, D., Mihăiţă, A.-S., Ou, Y., Grzybowska, H., Li, M. (2024). Origin–destination matrix estimation for public transport: A multi-modal weighted graph approach. \textit{Transportation Research Part C}, 140, 104694. https://doi.org/10.1016/j.trc.2024.104694
    \bibitem{fielbaum2016} Andrés Fielbaum, Sergio Jara-Diaz, Antonio Gschwender (2016). Optimal public transport networks for a general urban structure. \textit{Transportation Research Part B}, 94, 298.
    https://doi.org/10.1016/j.trb.2016.10.003
    \bibitem{ki2023} K.I. Cervantes-Sanmiguel, M.V. Chavez-Hernandez, O.J. Ibarra-Rojas. Analyzing the trade-off between minimizing travel times and
    reducing monetary costs for users in the transit network design. \textit{Transportation Research Part B}, 173, 142.
    https://doi.org/10.1016/j.trb.2023.04.009
    \bibitem{diestel2017} Diestel, R. (2017). \textit{Graph Theory}. Springer.
    \bibitem{balle1998} Bellá Bollobás (1998). \textit{Modern Graph Theory}. Springer.
    \bibitem{gobjalisco2024} Secretaría de Transporte del estado de Jalisco. \textit{Actualización de Rutas de Transporte Público Colectivo y Masivo en la Zona Metropolitana de Guadalajara}. Recuperado el 19 de septiembre del 2024 de \url{https://datos.jalisco.gob.mx/dataset/actualizacion-de-rutas-de-transporte-publico-colectivo-y-masivo-en-la-zona-metropolitana-de}.
    \bibitem{inegi2024} Instituto Nacional de Estadística y Geografía (INEGI).  \textit{Estadística de transporte urbano de pasajeros, enero de 2024: Indicadores económicos de coyuntura}. Recuperado el 1 de noviembre del 2024 de \url{https://inegi.org.mx/contenidos/programas/transporteurbano/doc/ETUP2024_03.pdf}.
    \bibitem{iieg2024} Instituto de Información Estadística y Geográfica de Jalisco. \textit{Usuarios en el Sistema de Transporte Urbano de Pasajeros de Guadalajara}. Recuperado el 2 de noviembre del 2024 de \url{https://iieg.gob.mx/ns/?page_id=25290}.

\end{thebibliography}

